\documentclass[11pt]{article}

    \usepackage[breakable]{tcolorbox}
    \usepackage{parskip} % Stop auto-indenting (to mimic markdown behaviour)
    

    % Basic figure setup, for now with no caption control since it's done
    % automatically by Pandoc (which extracts ![](path) syntax from Markdown).
    \usepackage{graphicx}
    % Keep aspect ratio if custom image width or height is specified
    \setkeys{Gin}{keepaspectratio}
    % Maintain compatibility with old templates. Remove in nbconvert 6.0
    \let\Oldincludegraphics\includegraphics
    % Ensure that by default, figures have no caption (until we provide a
    % proper Figure object with a Caption API and a way to capture that
    % in the conversion process - todo).
    \usepackage{caption}
    \DeclareCaptionFormat{nocaption}{}
    \captionsetup{format=nocaption,aboveskip=0pt,belowskip=0pt}

    \usepackage{float}
    \floatplacement{figure}{H} % forces figures to be placed at the correct location
    \usepackage{xcolor} % Allow colors to be defined
    \usepackage{enumerate} % Needed for markdown enumerations to work
    \usepackage{geometry} % Used to adjust the document margins
    \usepackage{amsmath} % Equations
    \usepackage{amssymb} % Equations
    \usepackage{textcomp} % defines textquotesingle
    % Hack from http://tex.stackexchange.com/a/47451/13684:
    \AtBeginDocument{%
        \def\PYZsq{\textquotesingle}% Upright quotes in Pygmentized code
    }


    \input{../../macros.tex} 
    \usepackage{upquote} % Upright quotes for verbatim code
    \usepackage{eurosym} % defines \euro

    \usepackage{iftex}
    \ifPDFTeX
        \usepackage[T1]{fontenc}
        \IfFileExists{alphabeta.sty}{
              \usepackage{alphabeta}
          }{
              \usepackage[mathletters]{ucs}
              \usepackage[utf8x]{inputenc}
          }
    \else
        \usepackage{fontspec}
        \usepackage{unicode-math}
    \fi

    \usepackage{fancyvrb} % verbatim replacement that allows latex
    \usepackage{grffile} % extends the file name processing of package graphics
                         % to support a larger range
    \makeatletter % fix for old versions of grffile with XeLaTeX
    \@ifpackagelater{grffile}{2019/11/01}
    {
      % Do nothing on new versions
    }
    {
      \def\Gread@@xetex#1{%
        \IfFileExists{"\Gin@base".bb}%
        {\Gread@eps{\Gin@base.bb}}%
        {\Gread@@xetex@aux#1}%
      }
    }
    \makeatother
    \usepackage[Export]{adjustbox} % Used to constrain images to a maximum size
    \adjustboxset{max size={0.9\linewidth}{0.9\paperheight}}

    % The hyperref package gives us a pdf with properly built
    % internal navigation ('pdf bookmarks' for the table of contents,
    % internal cross-reference links, web links for URLs, etc.)
    \usepackage{hyperref}
    % The default LaTeX title has an obnoxious amount of whitespace. By default,
    % titling removes some of it. It also provides customization options.
    \usepackage{titling}
    \usepackage{longtable} % longtable support required by pandoc >1.10
    \usepackage{booktabs}  % table support for pandoc > 1.12.2
    \usepackage{array}     % table support for pandoc >= 2.11.3
    \usepackage{calc}      % table minipage width calculation for pandoc >= 2.11.1
    \usepackage[inline]{enumitem} % IRkernel/repr support (it uses the enumerate* environment)
    \usepackage[normalem]{ulem} % ulem is needed to support strikethroughs (\sout)
                                % normalem makes italics be italics, not underlines
    \usepackage{soul}      % strikethrough (\st) support for pandoc >= 3.0.0
    \usepackage{mathrsfs}

    
    % Colors for the hyperref package
    \definecolor{urlcolor}{rgb}{0,.145,.698}
    \definecolor{linkcolor}{rgb}{.71,0.21,0.01}
    \definecolor{citecolor}{rgb}{.12,.54,.11}

    % ANSI colors
    \definecolor{ansi-black}{HTML}{3E424D}
    \definecolor{ansi-black-intense}{HTML}{282C36}
    \definecolor{ansi-red}{HTML}{E75C58}
    \definecolor{ansi-red-intense}{HTML}{B22B31}
    \definecolor{ansi-green}{HTML}{00A250}
    \definecolor{ansi-green-intense}{HTML}{007427}
    \definecolor{ansi-yellow}{HTML}{DDB62B}
    \definecolor{ansi-yellow-intense}{HTML}{B27D12}
    \definecolor{ansi-blue}{HTML}{208FFB}
    \definecolor{ansi-blue-intense}{HTML}{0065CA}
    \definecolor{ansi-magenta}{HTML}{D160C4}
    \definecolor{ansi-magenta-intense}{HTML}{A03196}
    \definecolor{ansi-cyan}{HTML}{60C6C8}
    \definecolor{ansi-cyan-intense}{HTML}{258F8F}
    \definecolor{ansi-white}{HTML}{C5C1B4}
    \definecolor{ansi-white-intense}{HTML}{A1A6B2}
    \definecolor{ansi-default-inverse-fg}{HTML}{FFFFFF}
    \definecolor{ansi-default-inverse-bg}{HTML}{000000}

    % common color for the border for error outputs.
    \definecolor{outerrorbackground}{HTML}{FFDFDF}

    % commands and environments needed by pandoc snippets
    % extracted from the output of `pandoc -s`
    \providecommand{\tightlist}{%
      \setlength{\itemsep}{0pt}\setlength{\parskip}{0pt}}
    \DefineVerbatimEnvironment{Highlighting}{Verbatim}{commandchars=\\\{\}}
    % Add ',fontsize=\small' for more characters per line
    \newenvironment{Shaded}{}{}
    \newcommand{\KeywordTok}[1]{\textcolor[rgb]{0.00,0.44,0.13}{\textbf{{#1}}}}
    \newcommand{\DataTypeTok}[1]{\textcolor[rgb]{0.56,0.13,0.00}{{#1}}}
    \newcommand{\DecValTok}[1]{\textcolor[rgb]{0.25,0.63,0.44}{{#1}}}
    \newcommand{\BaseNTok}[1]{\textcolor[rgb]{0.25,0.63,0.44}{{#1}}}
    \newcommand{\FloatTok}[1]{\textcolor[rgb]{0.25,0.63,0.44}{{#1}}}
    \newcommand{\CharTok}[1]{\textcolor[rgb]{0.25,0.44,0.63}{{#1}}}
    \newcommand{\StringTok}[1]{\textcolor[rgb]{0.25,0.44,0.63}{{#1}}}
    \newcommand{\CommentTok}[1]{\textcolor[rgb]{0.38,0.63,0.69}{\textit{{#1}}}}
    \newcommand{\OtherTok}[1]{\textcolor[rgb]{0.00,0.44,0.13}{{#1}}}
    \newcommand{\AlertTok}[1]{\textcolor[rgb]{1.00,0.00,0.00}{\textbf{{#1}}}}
    \newcommand{\FunctionTok}[1]{\textcolor[rgb]{0.02,0.16,0.49}{{#1}}}
    \newcommand{\RegionMarkerTok}[1]{{#1}}
    \newcommand{\ErrorTok}[1]{\textcolor[rgb]{1.00,0.00,0.00}{\textbf{{#1}}}}
    \newcommand{\NormalTok}[1]{{#1}}

    % Additional commands for more recent versions of Pandoc
    \newcommand{\ConstantTok}[1]{\textcolor[rgb]{0.53,0.00,0.00}{{#1}}}
    \newcommand{\SpecialCharTok}[1]{\textcolor[rgb]{0.25,0.44,0.63}{{#1}}}
    \newcommand{\VerbatimStringTok}[1]{\textcolor[rgb]{0.25,0.44,0.63}{{#1}}}
    \newcommand{\SpecialStringTok}[1]{\textcolor[rgb]{0.73,0.40,0.53}{{#1}}}
    \newcommand{\ImportTok}[1]{{#1}}
    \newcommand{\DocumentationTok}[1]{\textcolor[rgb]{0.73,0.13,0.13}{\textit{{#1}}}}
    \newcommand{\AnnotationTok}[1]{\textcolor[rgb]{0.38,0.63,0.69}{\textbf{\textit{{#1}}}}}
    \newcommand{\CommentVarTok}[1]{\textcolor[rgb]{0.38,0.63,0.69}{\textbf{\textit{{#1}}}}}
    \newcommand{\VariableTok}[1]{\textcolor[rgb]{0.10,0.09,0.49}{{#1}}}
    \newcommand{\ControlFlowTok}[1]{\textcolor[rgb]{0.00,0.44,0.13}{\textbf{{#1}}}}
    \newcommand{\OperatorTok}[1]{\textcolor[rgb]{0.40,0.40,0.40}{{#1}}}
    \newcommand{\BuiltInTok}[1]{{#1}}
    \newcommand{\ExtensionTok}[1]{{#1}}
    \newcommand{\PreprocessorTok}[1]{\textcolor[rgb]{0.74,0.48,0.00}{{#1}}}
    \newcommand{\AttributeTok}[1]{\textcolor[rgb]{0.49,0.56,0.16}{{#1}}}
    \newcommand{\InformationTok}[1]{\textcolor[rgb]{0.38,0.63,0.69}{\textbf{\textit{{#1}}}}}
    \newcommand{\WarningTok}[1]{\textcolor[rgb]{0.38,0.63,0.69}{\textbf{\textit{{#1}}}}}
    \makeatletter
    \newsavebox\pandoc@box
    \newcommand*\pandocbounded[1]{%
      \sbox\pandoc@box{#1}%
      % scaling factors for width and height
      \Gscale@div\@tempa\textheight{\dimexpr\ht\pandoc@box+\dp\pandoc@box\relax}%
      \Gscale@div\@tempb\linewidth{\wd\pandoc@box}%
      % select the smaller of both
      \ifdim\@tempb\p@<\@tempa\p@
        \let\@tempa\@tempb
      \fi
      % scaling accordingly (\@tempa < 1)
      \ifdim\@tempa\p@<\p@
        \scalebox{\@tempa}{\usebox\pandoc@box}%
      % scaling not needed, use as it is
      \else
        \usebox{\pandoc@box}%
      \fi
    }
    \makeatother

    % Define a nice break command that doesn't care if a line doesn't already
    % exist.
    \def\br{\hspace*{\fill} \\* }
    % Math Jax compatibility definitions
    \def\gt{>}
    \def\lt{<}
    \let\Oldtex\TeX
    \let\Oldlatex\LaTeX
    \renewcommand{\TeX}{\textrm{\Oldtex}}
    \renewcommand{\LaTeX}{\textrm{\Oldlatex}}
    % Document parameters
    % Document title
    \title{Homework\_07}
    
    
    
    
    
    
    
% Pygments definitions
\makeatletter
\def\PY@reset{\let\PY@it=\relax \let\PY@bf=\relax%
    \let\PY@ul=\relax \let\PY@tc=\relax%
    \let\PY@bc=\relax \let\PY@ff=\relax}
\def\PY@tok#1{\csname PY@tok@#1\endcsname}
\def\PY@toks#1+{\ifx\relax#1\empty\else%
    \PY@tok{#1}\expandafter\PY@toks\fi}
\def\PY@do#1{\PY@bc{\PY@tc{\PY@ul{%
    \PY@it{\PY@bf{\PY@ff{#1}}}}}}}
\def\PY#1#2{\PY@reset\PY@toks#1+\relax+\PY@do{#2}}

\@namedef{PY@tok@w}{\def\PY@tc##1{\textcolor[rgb]{0.73,0.73,0.73}{##1}}}
\@namedef{PY@tok@c}{\let\PY@it=\textit\def\PY@tc##1{\textcolor[rgb]{0.24,0.48,0.48}{##1}}}
\@namedef{PY@tok@cp}{\def\PY@tc##1{\textcolor[rgb]{0.61,0.40,0.00}{##1}}}
\@namedef{PY@tok@k}{\let\PY@bf=\textbf\def\PY@tc##1{\textcolor[rgb]{0.00,0.50,0.00}{##1}}}
\@namedef{PY@tok@kp}{\def\PY@tc##1{\textcolor[rgb]{0.00,0.50,0.00}{##1}}}
\@namedef{PY@tok@kt}{\def\PY@tc##1{\textcolor[rgb]{0.69,0.00,0.25}{##1}}}
\@namedef{PY@tok@o}{\def\PY@tc##1{\textcolor[rgb]{0.40,0.40,0.40}{##1}}}
\@namedef{PY@tok@ow}{\let\PY@bf=\textbf\def\PY@tc##1{\textcolor[rgb]{0.67,0.13,1.00}{##1}}}
\@namedef{PY@tok@nb}{\def\PY@tc##1{\textcolor[rgb]{0.00,0.50,0.00}{##1}}}
\@namedef{PY@tok@nf}{\def\PY@tc##1{\textcolor[rgb]{0.00,0.00,1.00}{##1}}}
\@namedef{PY@tok@nc}{\let\PY@bf=\textbf\def\PY@tc##1{\textcolor[rgb]{0.00,0.00,1.00}{##1}}}
\@namedef{PY@tok@nn}{\let\PY@bf=\textbf\def\PY@tc##1{\textcolor[rgb]{0.00,0.00,1.00}{##1}}}
\@namedef{PY@tok@ne}{\let\PY@bf=\textbf\def\PY@tc##1{\textcolor[rgb]{0.80,0.25,0.22}{##1}}}
\@namedef{PY@tok@nv}{\def\PY@tc##1{\textcolor[rgb]{0.10,0.09,0.49}{##1}}}
\@namedef{PY@tok@no}{\def\PY@tc##1{\textcolor[rgb]{0.53,0.00,0.00}{##1}}}
\@namedef{PY@tok@nl}{\def\PY@tc##1{\textcolor[rgb]{0.46,0.46,0.00}{##1}}}
\@namedef{PY@tok@ni}{\let\PY@bf=\textbf\def\PY@tc##1{\textcolor[rgb]{0.44,0.44,0.44}{##1}}}
\@namedef{PY@tok@na}{\def\PY@tc##1{\textcolor[rgb]{0.41,0.47,0.13}{##1}}}
\@namedef{PY@tok@nt}{\let\PY@bf=\textbf\def\PY@tc##1{\textcolor[rgb]{0.00,0.50,0.00}{##1}}}
\@namedef{PY@tok@nd}{\def\PY@tc##1{\textcolor[rgb]{0.67,0.13,1.00}{##1}}}
\@namedef{PY@tok@s}{\def\PY@tc##1{\textcolor[rgb]{0.73,0.13,0.13}{##1}}}
\@namedef{PY@tok@sd}{\let\PY@it=\textit\def\PY@tc##1{\textcolor[rgb]{0.73,0.13,0.13}{##1}}}
\@namedef{PY@tok@si}{\let\PY@bf=\textbf\def\PY@tc##1{\textcolor[rgb]{0.64,0.35,0.47}{##1}}}
\@namedef{PY@tok@se}{\let\PY@bf=\textbf\def\PY@tc##1{\textcolor[rgb]{0.67,0.36,0.12}{##1}}}
\@namedef{PY@tok@sr}{\def\PY@tc##1{\textcolor[rgb]{0.64,0.35,0.47}{##1}}}
\@namedef{PY@tok@ss}{\def\PY@tc##1{\textcolor[rgb]{0.10,0.09,0.49}{##1}}}
\@namedef{PY@tok@sx}{\def\PY@tc##1{\textcolor[rgb]{0.00,0.50,0.00}{##1}}}
\@namedef{PY@tok@m}{\def\PY@tc##1{\textcolor[rgb]{0.40,0.40,0.40}{##1}}}
\@namedef{PY@tok@gh}{\let\PY@bf=\textbf\def\PY@tc##1{\textcolor[rgb]{0.00,0.00,0.50}{##1}}}
\@namedef{PY@tok@gu}{\let\PY@bf=\textbf\def\PY@tc##1{\textcolor[rgb]{0.50,0.00,0.50}{##1}}}
\@namedef{PY@tok@gd}{\def\PY@tc##1{\textcolor[rgb]{0.63,0.00,0.00}{##1}}}
\@namedef{PY@tok@gi}{\def\PY@tc##1{\textcolor[rgb]{0.00,0.52,0.00}{##1}}}
\@namedef{PY@tok@gr}{\def\PY@tc##1{\textcolor[rgb]{0.89,0.00,0.00}{##1}}}
\@namedef{PY@tok@ge}{\let\PY@it=\textit}
\@namedef{PY@tok@gs}{\let\PY@bf=\textbf}
\@namedef{PY@tok@ges}{\let\PY@bf=\textbf\let\PY@it=\textit}
\@namedef{PY@tok@gp}{\let\PY@bf=\textbf\def\PY@tc##1{\textcolor[rgb]{0.00,0.00,0.50}{##1}}}
\@namedef{PY@tok@go}{\def\PY@tc##1{\textcolor[rgb]{0.44,0.44,0.44}{##1}}}
\@namedef{PY@tok@gt}{\def\PY@tc##1{\textcolor[rgb]{0.00,0.27,0.87}{##1}}}
\@namedef{PY@tok@err}{\def\PY@bc##1{{\setlength{\fboxsep}{\string -\fboxrule}\fcolorbox[rgb]{1.00,0.00,0.00}{1,1,1}{\strut ##1}}}}
\@namedef{PY@tok@kc}{\let\PY@bf=\textbf\def\PY@tc##1{\textcolor[rgb]{0.00,0.50,0.00}{##1}}}
\@namedef{PY@tok@kd}{\let\PY@bf=\textbf\def\PY@tc##1{\textcolor[rgb]{0.00,0.50,0.00}{##1}}}
\@namedef{PY@tok@kn}{\let\PY@bf=\textbf\def\PY@tc##1{\textcolor[rgb]{0.00,0.50,0.00}{##1}}}
\@namedef{PY@tok@kr}{\let\PY@bf=\textbf\def\PY@tc##1{\textcolor[rgb]{0.00,0.50,0.00}{##1}}}
\@namedef{PY@tok@bp}{\def\PY@tc##1{\textcolor[rgb]{0.00,0.50,0.00}{##1}}}
\@namedef{PY@tok@fm}{\def\PY@tc##1{\textcolor[rgb]{0.00,0.00,1.00}{##1}}}
\@namedef{PY@tok@vc}{\def\PY@tc##1{\textcolor[rgb]{0.10,0.09,0.49}{##1}}}
\@namedef{PY@tok@vg}{\def\PY@tc##1{\textcolor[rgb]{0.10,0.09,0.49}{##1}}}
\@namedef{PY@tok@vi}{\def\PY@tc##1{\textcolor[rgb]{0.10,0.09,0.49}{##1}}}
\@namedef{PY@tok@vm}{\def\PY@tc##1{\textcolor[rgb]{0.10,0.09,0.49}{##1}}}
\@namedef{PY@tok@sa}{\def\PY@tc##1{\textcolor[rgb]{0.73,0.13,0.13}{##1}}}
\@namedef{PY@tok@sb}{\def\PY@tc##1{\textcolor[rgb]{0.73,0.13,0.13}{##1}}}
\@namedef{PY@tok@sc}{\def\PY@tc##1{\textcolor[rgb]{0.73,0.13,0.13}{##1}}}
\@namedef{PY@tok@dl}{\def\PY@tc##1{\textcolor[rgb]{0.73,0.13,0.13}{##1}}}
\@namedef{PY@tok@s2}{\def\PY@tc##1{\textcolor[rgb]{0.73,0.13,0.13}{##1}}}
\@namedef{PY@tok@sh}{\def\PY@tc##1{\textcolor[rgb]{0.73,0.13,0.13}{##1}}}
\@namedef{PY@tok@s1}{\def\PY@tc##1{\textcolor[rgb]{0.73,0.13,0.13}{##1}}}
\@namedef{PY@tok@mb}{\def\PY@tc##1{\textcolor[rgb]{0.40,0.40,0.40}{##1}}}
\@namedef{PY@tok@mf}{\def\PY@tc##1{\textcolor[rgb]{0.40,0.40,0.40}{##1}}}
\@namedef{PY@tok@mh}{\def\PY@tc##1{\textcolor[rgb]{0.40,0.40,0.40}{##1}}}
\@namedef{PY@tok@mi}{\def\PY@tc##1{\textcolor[rgb]{0.40,0.40,0.40}{##1}}}
\@namedef{PY@tok@il}{\def\PY@tc##1{\textcolor[rgb]{0.40,0.40,0.40}{##1}}}
\@namedef{PY@tok@mo}{\def\PY@tc##1{\textcolor[rgb]{0.40,0.40,0.40}{##1}}}
\@namedef{PY@tok@ch}{\let\PY@it=\textit\def\PY@tc##1{\textcolor[rgb]{0.24,0.48,0.48}{##1}}}
\@namedef{PY@tok@cm}{\let\PY@it=\textit\def\PY@tc##1{\textcolor[rgb]{0.24,0.48,0.48}{##1}}}
\@namedef{PY@tok@cpf}{\let\PY@it=\textit\def\PY@tc##1{\textcolor[rgb]{0.24,0.48,0.48}{##1}}}
\@namedef{PY@tok@c1}{\let\PY@it=\textit\def\PY@tc##1{\textcolor[rgb]{0.24,0.48,0.48}{##1}}}
\@namedef{PY@tok@cs}{\let\PY@it=\textit\def\PY@tc##1{\textcolor[rgb]{0.24,0.48,0.48}{##1}}}

\def\PYZbs{\char`\\}
\def\PYZus{\char`\_}
\def\PYZob{\char`\{}
\def\PYZcb{\char`\}}
\def\PYZca{\char`\^}
\def\PYZam{\char`\&}
\def\PYZlt{\char`\<}
\def\PYZgt{\char`\>}
\def\PYZsh{\char`\#}
\def\PYZpc{\char`\%}
\def\PYZdl{\char`\$}
\def\PYZhy{\char`\-}
\def\PYZsq{\char`\'}
\def\PYZdq{\char`\"}
\def\PYZti{\char`\~}
% for compatibility with earlier versions
\def\PYZat{@}
\def\PYZlb{[}
\def\PYZrb{]}
\makeatother


    % For linebreaks inside Verbatim environment from package fancyvrb.
    \makeatletter
        \newbox\Wrappedcontinuationbox
        \newbox\Wrappedvisiblespacebox
        \newcommand*\Wrappedvisiblespace {\textcolor{red}{\textvisiblespace}}
        \newcommand*\Wrappedcontinuationsymbol {\textcolor{red}{\llap{\tiny$\m@th\hookrightarrow$}}}
        \newcommand*\Wrappedcontinuationindent {3ex }
        \newcommand*\Wrappedafterbreak {\kern\Wrappedcontinuationindent\copy\Wrappedcontinuationbox}
        % Take advantage of the already applied Pygments mark-up to insert
        % potential linebreaks for TeX processing.
        %        {, <, #, %, $, ' and ": go to next line.
        %        _, }, ^, &, >, - and ~: stay at end of broken line.
        % Use of \textquotesingle for straight quote.
        \newcommand*\Wrappedbreaksatspecials {%
            \def\PYGZus{\discretionary{\char`\_}{\Wrappedafterbreak}{\char`\_}}%
            \def\PYGZob{\discretionary{}{\Wrappedafterbreak\char`\{}{\char`\{}}%
            \def\PYGZcb{\discretionary{\char`\}}{\Wrappedafterbreak}{\char`\}}}%
            \def\PYGZca{\discretionary{\char`\^}{\Wrappedafterbreak}{\char`\^}}%
            \def\PYGZam{\discretionary{\char`\&}{\Wrappedafterbreak}{\char`\&}}%
            \def\PYGZlt{\discretionary{}{\Wrappedafterbreak\char`\<}{\char`\<}}%
            \def\PYGZgt{\discretionary{\char`\>}{\Wrappedafterbreak}{\char`\>}}%
            \def\PYGZsh{\discretionary{}{\Wrappedafterbreak\char`\#}{\char`\#}}%
            \def\PYGZpc{\discretionary{}{\Wrappedafterbreak\char`\%}{\char`\%}}%
            \def\PYGZdl{\discretionary{}{\Wrappedafterbreak\char`\$}{\char`\$}}%
            \def\PYGZhy{\discretionary{\char`\-}{\Wrappedafterbreak}{\char`\-}}%
            \def\PYGZsq{\discretionary{}{\Wrappedafterbreak\textquotesingle}{\textquotesingle}}%
            \def\PYGZdq{\discretionary{}{\Wrappedafterbreak\char`\"}{\char`\"}}%
            \def\PYGZti{\discretionary{\char`\~}{\Wrappedafterbreak}{\char`\~}}%
        }
        % Some characters . , ; ? ! / are not pygmentized.
        % This macro makes them "active" and they will insert potential linebreaks
        \newcommand*\Wrappedbreaksatpunct {%
            \lccode`\~`\.\lowercase{\def~}{\discretionary{\hbox{\char`\.}}{\Wrappedafterbreak}{\hbox{\char`\.}}}%
            \lccode`\~`\,\lowercase{\def~}{\discretionary{\hbox{\char`\,}}{\Wrappedafterbreak}{\hbox{\char`\,}}}%
            \lccode`\~`\;\lowercase{\def~}{\discretionary{\hbox{\char`\;}}{\Wrappedafterbreak}{\hbox{\char`\;}}}%
            \lccode`\~`\:\lowercase{\def~}{\discretionary{\hbox{\char`\:}}{\Wrappedafterbreak}{\hbox{\char`\:}}}%
            \lccode`\~`\?\lowercase{\def~}{\discretionary{\hbox{\char`\?}}{\Wrappedafterbreak}{\hbox{\char`\?}}}%
            \lccode`\~`\!\lowercase{\def~}{\discretionary{\hbox{\char`\!}}{\Wrappedafterbreak}{\hbox{\char`\!}}}%
            \lccode`\~`\/\lowercase{\def~}{\discretionary{\hbox{\char`\/}}{\Wrappedafterbreak}{\hbox{\char`\/}}}%
            \catcode`\.\active
            \catcode`\,\active
            \catcode`\;\active
            \catcode`\:\active
            \catcode`\?\active
            \catcode`\!\active
            \catcode`\/\active
            \lccode`\~`\~
        }
    \makeatother

    \let\OriginalVerbatim=\Verbatim
    \makeatletter
    \renewcommand{\Verbatim}[1][1]{%
        %\parskip\z@skip
        \sbox\Wrappedcontinuationbox {\Wrappedcontinuationsymbol}%
        \sbox\Wrappedvisiblespacebox {\FV@SetupFont\Wrappedvisiblespace}%
        \def\FancyVerbFormatLine ##1{\hsize\linewidth
            \vtop{\raggedright\hyphenpenalty\z@\exhyphenpenalty\z@
                \doublehyphendemerits\z@\finalhyphendemerits\z@
                \strut ##1\strut}%
        }%
        % If the linebreak is at a space, the latter will be displayed as visible
        % space at end of first line, and a continuation symbol starts next line.
        % Stretch/shrink are however usually zero for typewriter font.
        \def\FV@Space {%
            \nobreak\hskip\z@ plus\fontdimen3\font minus\fontdimen4\font
            \discretionary{\copy\Wrappedvisiblespacebox}{\Wrappedafterbreak}
            {\kern\fontdimen2\font}%
        }%

        % Allow breaks at special characters using \PYG... macros.
        \Wrappedbreaksatspecials
        % Breaks at punctuation characters . , ; ? ! and / need catcode=\active
        \OriginalVerbatim[#1,codes*=\Wrappedbreaksatpunct]%
    }
    \makeatother

    % Exact colors from NB
    \definecolor{incolor}{HTML}{303F9F}
    \definecolor{outcolor}{HTML}{D84315}
    \definecolor{cellborder}{HTML}{CFCFCF}
    \definecolor{cellbackground}{HTML}{F7F7F7}

    % prompt
    \makeatletter
    \newcommand{\boxspacing}{\kern\kvtcb@left@rule\kern\kvtcb@boxsep}
    \makeatother
    \newcommand{\prompt}[4]{
        {\ttfamily\llap{{\color{#2}[#3]:\hspace{3pt}#4}}\vspace{-\baselineskip}}
    }
    

    
    % Prevent overflowing lines due to hard-to-break entities
    \sloppy
    % Setup hyperref package
    \hypersetup{
      breaklinks=true,  % so long urls are correctly broken across lines
      colorlinks=true,
      urlcolor=urlcolor,
      linkcolor=linkcolor,
      citecolor=citecolor,
      }
    % Slightly bigger margins than the latex defaults
    
    \geometry{verbose,tmargin=1in,bmargin=1in,lmargin=1in,rmargin=1in}
    
    

\begin{document}
    
    \maketitle
    
    

    
    Probability for Data Science

UC Berkeley, Spring 2026

Ani Adhikari

CC BY-NC-SA 4.0

This content is protected and may not be shared, uploaded, or
distributed.

    \section{Homework 7 (Due March 9th at 5
PM)}\label{homework-7-due-march-9th-at-5-pm}

    \begin{tcolorbox}[breakable, size=fbox, boxrule=1pt, pad at break*=1mm,colback=cellbackground, colframe=cellborder]
\prompt{In}{incolor}{1}{\boxspacing}
\begin{Verbatim}[commandchars=\\\{\}]
\PY{k+kn}{import}\PY{+w}{ }\PY{n+nn}{warnings}
\PY{n}{warnings}\PY{o}{.}\PY{n}{filterwarnings}\PY{p}{(}\PY{l+s+s1}{\PYZsq{}}\PY{l+s+s1}{ignore}\PY{l+s+s1}{\PYZsq{}}\PY{p}{)}

\PY{k+kn}{from}\PY{+w}{ }\PY{n+nn}{prob140}\PY{+w}{ }\PY{k+kn}{import} \PY{o}{*}
\PY{k+kn}{from}\PY{+w}{ }\PY{n+nn}{datascience}\PY{+w}{ }\PY{k+kn}{import} \PY{o}{*}
\PY{k+kn}{import}\PY{+w}{ }\PY{n+nn}{numpy}\PY{+w}{ }\PY{k}{as}\PY{+w}{ }\PY{n+nn}{np}
\PY{k+kn}{from}\PY{+w}{ }\PY{n+nn}{scipy}\PY{+w}{ }\PY{k+kn}{import} \PY{n}{stats}

\PY{k+kn}{import}\PY{+w}{ }\PY{n+nn}{matplotlib}\PY{n+nn}{.}\PY{n+nn}{pyplot}\PY{+w}{ }\PY{k}{as}\PY{+w}{ }\PY{n+nn}{plt}
\PY{o}{\PYZpc{}}\PY{k}{matplotlib} inline
\PY{k+kn}{import}\PY{+w}{ }\PY{n+nn}{matplotlib}
\PY{n}{matplotlib}\PY{o}{.}\PY{n}{style}\PY{o}{.}\PY{n}{use}\PY{p}{(}\PY{l+s+s1}{\PYZsq{}}\PY{l+s+s1}{fivethirtyeight}\PY{l+s+s1}{\PYZsq{}}\PY{p}{)}
\end{Verbatim}
\end{tcolorbox}

    \subsubsection{Instructions}\label{instructions}

\paragraph{\texorpdfstring{\textbf{Questions 1, 2, and 3 can be answered
based on Tuesday's lecture (and past content)
alone.}}{Questions 1, 2, and 3 can be answered based on Tuesday's lecture (and past content) alone.}}\label{questions-1-2-and-3-can-be-answered-based-on-tuesdays-lecture-and-past-content-alone.}

Your homeworks will generally have two components: a written portion and
a portion that also involves code. Written work should be completed on
paper, and coding questions should be done in the notebook. Start the
work for the written portions of each section on a new page. You are
welcome to \(\LaTeX\) your answers to the written portions, but staff
will not be able to assist you with \(\LaTeX\) related issues.

It is your responsibility to ensure that both components of the homework
are submitted completely and properly to Gradescope. \textbf{Make sure
to assign each page of your pdf to the correct question. Refer to the
bottom of the notebook for submission instructions.}

    \subsubsection{How to Do Your Homework}\label{how-to-do-your-homework}

The point of homework is for you to try your hand at using what you've
learned in class. The steps to follow:

\begin{itemize}
\tightlist
\item
  Go to lecture and sections, and also go over the relevant text
  sections before starting on the homework. This will remind you what
  was covered in class, and the text will typically contain examples not
  covered in lecture. The weekly Study Guide will list what you should
  read.
\item
  Work on some of the practice problems before starting on the homework.
\item
  Attempt the homework problems by yourself with the text, section work,
  and practice materials all at hand. Sometimes the week's lab will help
  as well. The two steps above will help this step go faster and be more
  fruitful.
\item
  At this point, seek help if you need it. Don't ask how to do the
  problem --- ask how to get started, or for a nudge to get you past
  where you are stuck. Always say what you have already tried. That
  helps us help you more effectively.
\item
  For a good measure of your understanding, keep track of the fraction
  of the homework you can do by yourself or with minimal help. It's a
  better measure than your homework score, and only you can measure it.
\end{itemize}

    \subsubsection{Rules for Homework}\label{rules-for-homework}

\begin{itemize}
\tightlist
\item
  Every answer should contain a calculation or reasoning. For example, a
  calculation such as \((1/3)(0.8) + (2/3)(0.7)\) or
  \texttt{sum({[}(1/3)*0.8,\ (2/3)*0.7{]})}is fine without further
  explanation or simplification. If we want you to simplify, we'll ask
  you to. But just \({5 \choose 2}\) by itself is not fine; write ``we
  want any 2 out of the 5 frogs and they can appear in any order'' or
  whatever reasoning you used. Reasoning can be brief and abbreviated,
  e.g.~``product rule'' or ``not mutually exclusive.''
\item
  You may consult others (see ``How to Do Your Homework'' above) but you
  must write up your own answers using your own words, notation, and
  sequence of steps.
\item
  We'll be using Gradescope. You must submit the homework according to
  the instructions at the end of homework set.
\end{itemize}

    \subsection{\texorpdfstring{\textbf{We will not grade assignments which
do not have pages selected for each
question.}}{We will not grade assignments which do not have pages selected for each question.}}\label{we-will-not-grade-assignments-which-do-not-have-pages-selected-for-each-question.}

    \newpage

    \subsection{1. Correlation}\label{correlation}

The \emph{correlation coefficient} between random variables \(X\) and
\(Y\) is defined as

\[
r(X, Y) ~ = \frac{Cov(X, Y)}{SD(X)SD(Y)}
\]

It is called the correlation, for short. The definition explains why
\(X\) and \(Y\) are called \emph{uncorrelated} if \(Cov(X, Y) = 0\).

\textbf{a)} Let \(X^*\) be \(X\) in standard units and let \(Y^*\) be
\(Y\) in standard units. Check that

\[
r(X, Y) = E(X^*Y^*)
\]


\begin{align*}
	r(X,Y) &= \frac{\Cov(X,Y)}{SD(X)SD(Y)} \\
	&=  \frac{E[(X-\mu_{X}) \cdot (Y-\mu _{Y})]}{SD(X)SD(Y)} \\
	&= E\left[ \frac{1}{SD(X)SD(Y)} (X-\mu _{X}) (Y-\mu _{Y})\right], \quad \text{by linearity of expectation}\\ 
	&= E\left[ \frac{X-\mu _{X}}{SD(X)} \cdot \frac{Y-\mu _{Y}}{SD(Y)} \right] \\
	&= E[X^{*}Y^{*}] \\
\end{align*}



This is the random variable version of the Data 8 definition of the
correlation between two data variables: convert each variable to
standard units; multiply each pair; take the mean of the products.

\textbf{b)} Use the fact that \((X^* + Y^*)^2\) and \((X^* - Y^*)^2\)
are non-negative random variables to show that \(-1 \le r(X, Y) \le 1\).

{[}First find the numerical values of \(E(X^*)\) and
\(E\left({X^*}^2\right)\). Then find \(E\left((X^* + Y^*)^2\right)\).{]}

\begin{align*}
	E[X^{*}] &= E\left[\frac{X-\mu _{X}}{\sigma_{X}}\right] \\ 
		 &= \frac{1}{\sigma_{X}} \cdot E[X-\mu_{X}] \\
		 &= \frac{1}{\sigma_{X}} \cdot E[X] - \mu_{X}  \\
		 &= \frac{1}{\sigma_{X}} \cdot E[X] -E[X] \\
		 &= 0 \\
\end{align*}


\begin{align*}
	E[X^{*^{2}}] &= E\left[\frac{1}{\sigma^{2}} (X-\mu _{X})^{2}\right] \\
		     &= \frac{1}{\sigma^{2}}E\left[ X^{2} -2X\mu _{X} + \mu _{X}^{2}\right] \\
		     &= \frac{1}{\sigma^{2}}\left( E[X^{2}] - 2\mu _{x}E[X] + \mu ^{2}_{X} \right)   \\
		     &= \frac{1}{\sigma^{2}} (E[X^{2}] -\mu _{X}^{2}) \\
		     &= \frac{1}{\sigma^{2}} \cdot \Var(X) \\
		     &= \frac{\Var(X)}{\Var(X)} \\
		     &= 1 \\
\end{align*}

\begin{align*}
	E[(X^{*}+Y^{*})^{2}] &= E[X^{*^{2}} +2X^{*}Y^{*} + Y^{*^{2}} ]  \\
	&= E[X^{*^{2}}] + E[Y^{*^{2}}] + 2E[X^{*}Y^{*}] \\
	&= 1 + 1 + 2r(X,Y)  \\
	&= 2+2r(X,Y) \\
\end{align*}


\begin{align*}
	E[(X^{*} - Y^{*})^{2}] &= E[X^{*}] + E[Y^{*}] - 2E[X^{*}Y^{*}] \\
	&= 2 - 2r(X,Y) \\
\end{align*}

Since $(X^* + Y^*)^2$ and $(X^* - Y^*)^2$
are non-negative random variables, that means the expectation of each of those random variables also has to be positive. Therefore: 

\begin{align*}
	2+2r(X,Y) &\ge 0\\
	2r(X,Y) &\ge -2\\
	r(X,Y) &\ge -1\\
\end{align*}

\begin{align*}
	2-2r(X,Y) &\ge 0\\
	r(X,Y) &\ge 1
\end{align*}

Thus, $-1 \le r(X, Y) \le 1$


\textbf{c)} Show that if \(Y = aX+b\) where \(a \ne 0\), then
\(r(X, Y)\) is 1 or \(-1\) depending on whether the sign of \(a\) is
positive or negative.


\begin{align*}
	r(X,Y) &= \frac{1}{SD(X)SD(Y)} \cdot \Cov(X,Y) \\
	&= \frac{1}{SD(X)SD(aX+b)} \cdot \Cov(X, aX+b) \\
\end{align*}

\begin{align*}
	\Cov(X, aX+b) &= \Cov(X, aX) \\
	&= a\Cov(X, X) \\
	&= a\Var(X) \\
\end{align*}

\begin{align*}
	SD(Y) &= SD(aX+b) \\
	      &= \sqrt{\Var(aX+b)} \\
	      &= \sqrt{\Var(aX)}\\
	      &= \sqrt{a^{2}\Var(X)} \\
	      &= |a|\cdot SD(X) \\
\end{align*}

\begin{align*}
	r(X,Y) &= \frac{1}{|a|\Var(X)} \cdot a\Var(X) \\
	&= \frac{a}{|a|} \\
\end{align*}

Thus, the sign of $ r(X,Y) $ is $ -1 $ and $ 1 $ depending on a. 

\textbf{d)} Consider a sequence of i.i.d. Bernoulli \((p)\) trials. For
any positive integer \(k\) let \(X_k\) be the number of successes in
trials 1 through \(k\). \textbf{Use bilinearity} to find
\(Cov(X_n, X_{n+m})\) and hence find \(r(X_n, X_{n+m})\).


We can think of $ X_{n+m} $ as the number of successes from $ 1,2,3,\dots n $ and $ n+1,n+2,n+3,\dots, n+m $.

Since it is i.i.d. Bernoulli trials, the number of successes in trials  $ n+1, n+2, \dots, n+m $ is the same as $ 1,2,\dots m $, or $ X_{m} $, independent of $ X_{n} $  

\begin{align*}
	\Cov(X_{n}, X_{n+m}) &= \Cov(X_{n}, X_{n} + X_{m}) \\
	&= \Cov(X_{n}, X_{n}) + \Cov(X_{n}, X_{m}) \\
	&= \Var(X_{n}) + 0 \\
	&= \Var(\sum_{i=1}^{n}X_{i}) \\
	&= n\Var(X_{1}) \quad \text{by symmetry} \\
	&= n\cdot p(1-p) \\
\end{align*}







\textbf{e)} Fix \(n\) and find the limit of \(r(X_n, X_{n+m})\) as
\(m \to \infty\). Explain why the limit is consistent with intuition.

\begin{align*}
	\lim_{m\to \infty }r(X_{n}, X_{n+m}) &= \lim_{m\to\infty } \frac{\Cov(X_{n},X_{n+m})}{SD(X_{n})SD(X_{n+m})} \\
	&= \lim_{m\to\infty } \frac{n \cdot p(1-p)}{SD(X_{n}) SD(X_{n+m})} \\
	&= \lim_{m\to \infty } \frac{n \cdot p(1-p)}{\sqrt{np(1-p)} \cdot \sqrt{(n+m)p(1-p)}} \\
	&= 0 \\
\end{align*}

The numerator does not depend on $ m $, while the denominator grows to $ \infty  $ as $ m \to \infty  $. Therefore the limit is 0. This makes sense because as the correlation between the two variables will only depend on the first $ n $ trials, then account for the variance of two variables. Since the variance of $ X_{n+m} $ increases as there are more trials conducted, then the correlation will keep decreasing. Taking the limit gives 0. 







    \newpage

    \subsection*{2. Collecting Distinct
Values}\label{collecting-distinct-values}

In Homework 4 you found the expectation of each of the random variables
below. \textbf{Go back and see how you did that, and then use the same
ideas} to find the variance of each one.

For one part you will need the fact that the SD of a geometric \((p)\)
random variable is \(\frac{\sqrt{q}}{p}\) where \(q = 1-p\). We haven't
proved that as the algebra takes a bit of work. We will prove it later
in the course by conditioning.

\textbf{(a)} A die is rolled \(n\) times. Find the variance of number of
faces that \emph{do not} appear.

Let $ X $ be the number of faces that do not appear.

Let $ I_{n} = \begin{cases}
	1, &\text{if the side with $ n $ dots does not show up}\\
	0, &\text{otherwise}
\end{cases}, \quad I_{n} \sim \Bern((\frac{5}{6})^{n}) $

$ I_{n} $ is i.i.d. of all other indicators. 
\[
X = \sum_{n=1}^{6}I_{n}
\]

\begin{align*}
	\Var(X) &= \Var(\sum_{n=1}^{6}I_{n}) \\
&= \sum_{i=1}^{6}\sum_{j=1}^{6}\Cov(I_{i}, I_{j}) \\
&= \sum_{i=1}^{6} \Var(I_{i}) + \mathop{\sum\sum}_{1\le i\neq j \le n} \Cov(I_{i}, I_{j})\\
&= \sum_{i=1}^{6}\Var(I_{i}) + 0 \quad \text{by independence} \\
&= 6\Var(I_{1}) \\
&= 6p(1-p) \\
&= 6 \left[  \left( \frac{5}{6} \right)^{n} (1- \left( \frac{5}{6} \right)^{n})\right]  \\
\end{align*}




\textbf{(b)} Use your answer to (a) to find the variance of the number
of distinct faces that \emph{do} appear in \(n\) rolls of a die.

\begin{align*}
	\Var(6-X) &= \Var(X) \\
&= 6 \left[  \left( \frac{5}{6} \right)^{n} (1- \left( \frac{5}{6} \right)^{n})\right]  \\
\end{align*}



\textbf{(c)} Find the variance of the number of times you have to roll a
die till you have seen all of the faces.

Let $ Y $ be the number of times you have to roll a die till you seen all the faces. 

Let $ Y_{i} $ be the number of times you have to roll a die until you have seen the $ i $th new face after already seeing $ i-1 $ faces. Where $ Y_{i} \sim Geometric(\frac{6-(i-1)}{6}) $  

Then, 
\[
Y = \sum_{i=1}^{6}Y_{i} 
\]

And $ Y_{1} \perp Y_{2} \perp Y_{3} \perp \dots \perp Y_{6} $ 

\begin{align*}
	\Var(Y) &= \Var(\sum_{i=1}^{6}Y_{i}) \\
	&= \sum_{i=1}^{6}\Var(Y_{i}) + 0 \quad \text{by independence}\\
	&= \Var(Y_1) + \Var(Y_{2}) + \dots \Var(Y_{6})\\
	&= \frac{\sqrt{\frac{0}{6}}}{\frac{6}{6}} + \frac{\sqrt{\frac{1}{6}}}{\frac{5}{6}} + \frac{\sqrt{\frac{2}{6}}}{\frac{4}{6}} + \dots + \frac{\sqrt{\frac{5}{6}}}{\frac{1}{6}} \\
	&\approx 10.69  
\end{align*}


    \newpage

    \subsection*{3. The ``Sample Variance''}\label{the-sample-variance}

Let \(X_1, X_2, \ldots, X_n\) be i.i.d. (discrete or continuous), each
with mean \(\mu\) and SD \(\sigma\). Let
\(\bar{X} = \frac{1}{n}\sum_{i=1}^n X_i\) be the sample mean.

\textbf{(a)} Find \(E(\bar{X})\) and \(SD(\bar{X})\). It's fine to just
quote answers derived in class or in the textbook.


\begin{align*}
	E[\bar{X}] &= E[\frac{1}{n}\sum_{n=1}^{n}X_{i}] \\
	&= \frac{1}{n}E[\sum_{n=1}^{n}X_{i}]\\
	&= \frac{1}{n} n\cdot E[X_{1}] \\
	&= E[X_{1}] \\
	&= \mu  \\
\end{align*}

\begin{align*}
	\Var(\bar X) &= \sum_{i=1}^{n}\sum_{j=1}^{n}\Cov(X_{i},X_{j}) \\
	&= \sum_{i=1}^{n}\Var(X_{ia}) \\
	&= n\Var(X_1) \\
	&= n\sigma^{2} \\
\end{align*}

 


\textbf{(b)} For each \(i\), find \(Cov(X_i, \bar{X})\). {[}Plug in the
definition of \(\bar{X}\) and use bilinearity.{]}

\begin{align*}
	\Cov(X_{i}, \bar X) &= \Cov(X_{i}, \frac{1}{n}\sum_{n=1}^{n}X_{n}) \\
	&= \frac{1}{n}\Cov(X_{i}, \sum_{n=1}^{n}X_{n}) \\
	&= \frac{1}{n}\left( \Cov(X_{i}, X_{i}) + (n-1) \cdot 0\right), \quad \text{by independence} \\
	&= \frac{1}{n}\Cov(X_{i}, X_{i}) \\
	&= \frac{1}{n}\Var(X_{i}) \\
	&= \frac{\sigma^{2}}{n} \forall i \\
\end{align*}


\textbf{(c)} For each \(i\) in the range 1 through \(n\), define the
\emph{\(i\)th deviation in the sample} as \(D_i = X_i - \bar{X}\). Find
\(E(D_i)\) and \(Var(D_i)\). {[}Write the variance as \(Cov(D_i, D_i)\),
plug in the definition of \(D_i\), and use bilinearity.{]}

\begin{align*}
	E[D_{i}] &= E[X_{i} - \bar X] \\
	&= E[X_{i}] - E[\bar X] \\
	&= \mu  - \mu  \\
	&= 0 \\
\end{align*}

\begin{align*}
	\Var(D_{i}) &= \Cov(D_{i}, D_{i}) \\
	&= \Cov(X_{i}-\bar X, X_{i} - \bar X) \\
	&= \Cov(X_{i}, X_{i}) + \Cov(-\bar X, X_{i}) + \Cov(-\bar X, X_{i}) + \Cov(-\bar X, - \bar X) \\
	&= \Var(X_{i}) -2 \Cov(X_{i}, \bar X) + \Cov(\bar X, \bar X) \\
	&= \sigma^{2} -2 \frac{\sigma^{2}}{n} + \frac{1}{n}\sigma^{2} \\
	&= \sigma^{2}(1-\frac{1}{n}) \\
	&= \sigma^{2}(\frac{n-1}{n}) \\
\end{align*}


\textbf{(d)} Define the random variable \(\hat{\sigma}^2\) as \[
\hat{\sigma}^2 ~ = ~ \frac{1}{n} \sum_{i=1}^n D_i^2
\] Find \(E(\hat{\sigma}^2)\).

For this random variable, the notation \(\hat{\sigma}^2\) is pretty
standard in statistics. Just think of \(\hat{\sigma}^2\) as a symbol; it
doesn't help to start thinking about the random variable that is its
square root.


\begin{align*}
	E[\hat \sigma^{2}] &= E\left[\frac{1}{n} \sum_{i=1}^{n}D_{i}^{2}\right] \\
	&= \frac{1}{n}\sum_{i=1}^{n}E[D_{i}^{2}] \\ 
	&= \frac{1}{n} \cdot nE[D_1{^{2}}], \quad \text{by symmetry} \\
	&= E[D_1^{2}] \\
\end{align*}

\begin{align*}
	\Var(D_{i}) &= E[D_{i}^{2}] - E[D_{i}]^{2} \\
	&= E[D_{i}^{2}] - 0 \\
\end{align*}

\begin{align*}
	E[\hat \sigma^{2}] &= \Var(D_{i}) \\
	&= \sigma^{2} (\frac{n-1}{n}) \\
\end{align*}
\textbf{(e)} Use Part \textbf{d} to construct a random variable denoted
\(S^2\) that is an unbiased estimator of \(\sigma^2\). This random
variable \(S^2\) is called the \emph{sample variance} and is frequently
used in inference.

We need $ E[S^{2}] = \sigma^{2} $ 


\begin{align*}
	E[\hat \sigma^{2}] &= \sigma^{2} \frac{n-1}{n} \\
	\sigma^{2} &= \frac{n}{n-1} \cdot E[\hat \sigma^{2}]  \\
	&= E[\frac{n}{n-1} \cdot \hat \sigma^{2}] \\
\end{align*}

Therefore we can define $ E[S^{2}] = E\left[ \left( \frac{n}{n-1} \hat \sigma^{2}\right) \right] $ 
    \newpage

    \subsection{4. Poisson-Binomial
Distribution}\label{poisson-binomial-distribution}

For this exercise, please refer to the theory in
\href{https://data140.org/textbook/content/chapter-14/exact-distribution-of-a-sum/}{Section
14.1} and the code in
\href{https://data140.org/textbook/content/chapter-14/pgfs-in-numpy/}{Section
14.2}.

In Lab 3B you saw that a \emph{Poisson-binomial} random variable is a
sum of independent indicators that are not necessarily identically
distributed:

\(X = I_1 + I_2 + \cdots + I_n\) where \(I_j\) has the Bernoulli
\((p_j)\) distribution and \(I_1, I_2, \ldots, I_n\) are independent.

\textbf{(a)} What is the probability generating function of a Bernoulli
\((p)\) random variable? Provide a formula and then use the code cell
below to define a function \texttt{indicator\_pgf} that takes \(p\) as
its argument and returns the probability generating function of a
Bernoulli \((p)\) random variable as a \texttt{NumPy} polynomial. Use as
many lines as you need. The last line of the cell is there for you to
check that your function is working.

    \begin{tcolorbox}[breakable, size=fbox, boxrule=1pt, pad at break*=1mm,colback=cellbackground, colframe=cellborder]
\prompt{In}{incolor}{ }{\boxspacing}
\begin{Verbatim}[commandchars=\\\{\}]
\PY{c+c1}{\PYZsh{} Answer to a}

\PY{k}{def}\PY{+w}{ }\PY{n+nf}{indicator\PYZus{}pgf}\PY{p}{(}\PY{n}{p}\PY{p}{)}\PY{p}{:}
    \PY{o}{.}\PY{o}{.}\PY{o}{.}
    \PY{k}{return} \PY{o}{.}\PY{o}{.}\PY{o}{.}

\PY{n+nb}{print}\PY{p}{(}\PY{n}{indicator\PYZus{}pgf}\PY{p}{(}\PY{l+m+mf}{0.4}\PY{p}{)}\PY{p}{)}
\end{Verbatim}
\end{tcolorbox}

    \textbf{(b)} For \(j = 1, 2, \ldots, 20\), let \(p_j = 1/(j+1)\). Let
\(I_1, I_2, \ldots, I_{20}\) be independent indicators such that \(I_j\)
has the Bernoulli \((p_j)\) distribution, and let
\(X = I_1 + I_2 + \cdots + I_{20}\). Complete the code cell below so
that \texttt{pgf\_X} is the probability generating function of \(X\) as
a \texttt{NumPy} polynomial. Use as many lines as you need. The last two
lines are there for you to check that your polynomial has the correct
degree and that it is indeed a probability generating function.

    \begin{tcolorbox}[breakable, size=fbox, boxrule=1pt, pad at break*=1mm,colback=cellbackground, colframe=cellborder]
\prompt{In}{incolor}{ }{\boxspacing}
\begin{Verbatim}[commandchars=\\\{\}]
\PY{c+c1}{\PYZsh{} Answer to b}

\PY{o}{.}\PY{o}{.}\PY{o}{.}

\PY{n+nb}{print}\PY{p}{(}\PY{n}{pgf\PYZus{}X}\PY{p}{)}
\PY{n+nb}{sum}\PY{p}{(}\PY{n}{pgf\PYZus{}X}\PY{o}{.}\PY{n}{c}\PY{p}{)} \PY{c+c1}{\PYZsh{} sum of coefficients}
\end{Verbatim}
\end{tcolorbox}

    \textbf{(c)} Complete the cell below to plot the probability histogram
of \(X\). Do not add any more lines.

    \begin{tcolorbox}[breakable, size=fbox, boxrule=1pt, pad at break*=1mm,colback=cellbackground, colframe=cellborder]
\prompt{In}{incolor}{ }{\boxspacing}
\begin{Verbatim}[commandchars=\\\{\}]
\PY{c+c1}{\PYZsh{} Answer to c}

\PY{n}{vals\PYZus{}X} \PY{o}{=} \PY{o}{.}\PY{o}{.}\PY{o}{.}
\PY{n}{probs\PYZus{}X} \PY{o}{=} \PY{o}{.}\PY{o}{.}\PY{o}{.}
\PY{n}{dist\PYZus{}X} \PY{o}{=} \PY{n}{Table}\PY{p}{(}\PY{p}{)}\PY{o}{.}\PY{o}{.}\PY{o}{.}
\PY{n}{Plot}\PY{p}{(}\PY{n}{dist\PYZus{}X}\PY{p}{)}
\end{Verbatim}
\end{tcolorbox}

    \textbf{(d)} Complete the cell below to find the expectation, variance,
and SD of \(X\) using \texttt{p\_array}. Do not add any more lines. Then
run the cell below that to check your answers.

    \begin{tcolorbox}[breakable, size=fbox, boxrule=1pt, pad at break*=1mm,colback=cellbackground, colframe=cellborder]
\prompt{In}{incolor}{ }{\boxspacing}
\begin{Verbatim}[commandchars=\\\{\}]
\PY{c+c1}{\PYZsh{} Answer to d}

\PY{n}{p\PYZus{}array} \PY{o}{=} \PY{l+m+mi}{1}\PY{o}{/}\PY{n}{np}\PY{o}{.}\PY{n}{arange}\PY{p}{(}\PY{l+m+mi}{2}\PY{p}{,} \PY{l+m+mi}{22}\PY{p}{)}
\PY{n}{ev\PYZus{}X} \PY{o}{=} \PY{o}{.}\PY{o}{.}\PY{o}{.}
\PY{n}{var\PYZus{}X} \PY{o}{=} \PY{o}{.}\PY{o}{.}\PY{o}{.}
\PY{n}{sd\PYZus{}X} \PY{o}{=} \PY{o}{.}\PY{o}{.}\PY{o}{.}
\PY{n}{ev\PYZus{}X}\PY{p}{,} \PY{n}{var\PYZus{}X}\PY{p}{,} \PY{n}{sd\PYZus{}X}
\end{Verbatim}
\end{tcolorbox}

    \begin{tcolorbox}[breakable, size=fbox, boxrule=1pt, pad at break*=1mm,colback=cellbackground, colframe=cellborder]
\prompt{In}{incolor}{ }{\boxspacing}
\begin{Verbatim}[commandchars=\\\{\}]
\PY{n}{dist\PYZus{}X}\PY{o}{.}\PY{n}{ev}\PY{p}{(}\PY{p}{)}\PY{p}{,} \PY{n}{dist\PYZus{}X}\PY{o}{.}\PY{n}{var}\PY{p}{(}\PY{p}{)}\PY{p}{,} \PY{n}{dist\PYZus{}X}\PY{o}{.}\PY{n}{sd}\PY{p}{(}\PY{p}{)}
\end{Verbatim}
\end{tcolorbox}

    \textbf{(e)} Explain why the distribution of \(X\) cannot be Poisson.
Then show that the distribution of \(X\) is not binomial either, as
follows. If \(X\) were binomial, what would \(n\) have to be? Use that
and your answer to Part (d) to see what \(p\) would have to be. Use the
code cell below to find the variance of that binomial distribution, and
compare with your answer to Part (d).

    \begin{tcolorbox}[breakable, size=fbox, boxrule=1pt, pad at break*=1mm,colback=cellbackground, colframe=cellborder]
\prompt{In}{incolor}{ }{\boxspacing}
\begin{Verbatim}[commandchars=\\\{\}]
\PY{c+c1}{\PYZsh{} Answer to e}

\PY{n}{n} \PY{o}{=} \PY{o}{.}\PY{o}{.}\PY{o}{.}
\PY{n}{p} \PY{o}{=} \PY{o}{.}\PY{o}{.}\PY{o}{.}
\PY{n}{binomial\PYZus{}variance} \PY{o}{=} \PY{o}{.}\PY{o}{.}\PY{o}{.} 
\PY{n}{binomial\PYZus{}variance}
\end{Verbatim}
\end{tcolorbox}

    \newpage

    \subsection{5. Testing Hypotheses in the Gauss
Model}\label{testing-hypotheses-in-the-gauss-model}

The \href{https://en.wikipedia.org/wiki/Carl_Friedrich_Gauss}{Gauss}
model for measurement error says that repeated measurements
\(X_1, X_2, \ldots, X_n\) of the same quantity have the structure

\[
X_i = \mu + \epsilon_i, ~~~~~ 1 \leq i \leq n
\]

where \(\mu\) is an unknown constant called the \emph{true value} and
\(\epsilon_1, \epsilon_2, \ldots, \epsilon_n\) are random error terms
assumed to be i.i.d. with mean 0 and variance \(\sigma^2\).

From a practical perspective, the true value \(\mu\) comes from the
quantity being measured (for example, the true weight of an object). The
error terms come from the measuring process (for example, from the
balance being used for weighing). Thus \(\sigma\) is sometimes known
because of extensive experience with the measuring process (for example,
having used the same balance to weigh many different objects).

So assume that the Gauss model holds with \(\sigma = 1\), and let
\(n = 100\). Suppose a data scientist wants to test the following
hypotheses:

\begin{itemize}
\tightlist
\item
  Null hypothesis \(H_0\): \(\mu = 20\)
\item
  Alternative hypothesis \(H_A\): \(\mu \neq 20\)
\end{itemize}

Suppose the data scientist wants to use the average measurement
\(\bar{X}\) as the test statistic and reject the null hypothesis if
\(\vert \bar{X} - 20 \vert > 0.175\).

\textbf{(a)} Rewrite the decision rule by filling in the blanks with
numbers:
\(\vert \bar{X} - 20 \vert > 0.175 \iff \bar{X} < \underline{~~~~~~~~} \text{ or } \bar{X} > \underline{~~~~~~~~}\)

\textbf{(b) Level:} Find the approximate distribution of the test
statistic \(\bar{X}\) under \(H_0\), and use this distribution to find
the approximate probability that the test rejects the null hypothesis if
the null hypothesis is true. This probability is called the \emph{level}
of the test. In Data 8 we called it the cutoff for the p-value.

\textbf{Please write out your answer in math notation.} Then use the
code cell below for scratch work. Remember that
\texttt{stats.norm.cdf(x,\ mean,\ SD)} evaluates to the cdf of the
normal \((\text{mean, } \text{SD}^2)\) distribution at the point \(x\).
The necessary modules have been imported at the top of this notebook.

\textbf{(c) Power:} Suppose that in fact \(\mu = 20.5\) though the data
scientist doesn't know this and is still performing the same test as
above. Find the approximate distribution of the test statistic
\(\bar{X}\) under the condition \(\mu = 20.5\), and use this
distribution to find the approximate probability that the test rejects
the null hypothesis if \(\mu = 20.5\). This probability is called the
\emph{power of the test against the fixed alternative \(\mu = 20.5\)}.

\textbf{Please write out your answer in math notation.} Then use the
code cell below for scratch work.

    \begin{tcolorbox}[breakable, size=fbox, boxrule=1pt, pad at break*=1mm,colback=cellbackground, colframe=cellborder]
\prompt{In}{incolor}{ }{\boxspacing}
\begin{Verbatim}[commandchars=\\\{\}]
\PY{c+c1}{\PYZsh{} Scratch work for b and c}
\end{Verbatim}
\end{tcolorbox}

    \textbf{(d)} Complete the code cell below to plot the graph of the power
of the test under the fixed alternative \(\mu = \mu_A\) for \(\mu_A\) in
the range \texttt{true\_mu} below. Do not add any more lines.

Computational note: First study the code below and compare with the
output of the cell.

    \begin{tcolorbox}[breakable, size=fbox, boxrule=1pt, pad at break*=1mm,colback=cellbackground, colframe=cellborder]
\prompt{In}{incolor}{ }{\boxspacing}
\begin{Verbatim}[commandchars=\\\{\}]
\PY{n}{mu\PYZus{}list} \PY{o}{=} \PY{p}{[}\PY{l+m+mi}{10}\PY{p}{,} \PY{l+m+mi}{15}\PY{p}{,} \PY{l+m+mi}{20}\PY{p}{]}  \PY{c+c1}{\PYZsh{} It\PYZsq{}s also fine for this to be an array.}

\PY{c+c1}{\PYZsh{} array of P(X\PYZus{}i \PYZlt{} 12)}
\PY{c+c1}{\PYZsh{} for X\PYZus{}i normal with mean = ith element of mu\PYZus{}list}
\PY{c+c1}{\PYZsh{} and SD = 8}
\PY{n}{stats}\PY{o}{.}\PY{n}{norm}\PY{o}{.}\PY{n}{cdf}\PY{p}{(}\PY{l+m+mi}{12}\PY{p}{,} \PY{n}{mu\PYZus{}list}\PY{p}{,} \PY{l+m+mi}{8}\PY{p}{)}
\end{Verbatim}
\end{tcolorbox}

    \begin{tcolorbox}[breakable, size=fbox, boxrule=1pt, pad at break*=1mm,colback=cellbackground, colframe=cellborder]
\prompt{In}{incolor}{ }{\boxspacing}
\begin{Verbatim}[commandchars=\\\{\}]
\PY{c+c1}{\PYZsh{} Answer to d}

\PY{n}{true\PYZus{}mu} \PY{o}{=} \PY{n}{np}\PY{o}{.}\PY{n}{arange}\PY{p}{(}\PY{l+m+mi}{19}\PY{p}{,} \PY{l+m+mi}{21}\PY{p}{,} \PY{l+m+mf}{0.01}\PY{p}{)}
\PY{n}{power} \PY{o}{=} \PY{o}{.}\PY{o}{.}\PY{o}{.}

\PY{n}{plt}\PY{o}{.}\PY{n}{plot}\PY{p}{(}\PY{n}{true\PYZus{}mu}\PY{p}{,} \PY{n}{power}\PY{p}{,} \PY{n}{color}\PY{o}{=}\PY{l+s+s1}{\PYZsq{}}\PY{l+s+s1}{darkblue}\PY{l+s+s1}{\PYZsq{}}\PY{p}{,} \PY{n}{lw}\PY{o}{=}\PY{l+m+mi}{2}\PY{p}{)}
\PY{n}{plt}\PY{o}{.}\PY{n}{xlabel}\PY{p}{(}\PY{l+s+sa}{r}\PY{l+s+s1}{\PYZsq{}}\PY{l+s+s1}{True value of \PYZdl{}}\PY{l+s+s1}{\PYZbs{}}\PY{l+s+s1}{mu\PYZdl{}}\PY{l+s+s1}{\PYZsq{}}\PY{p}{)}
\PY{n}{plt}\PY{o}{.}\PY{n}{title}\PY{p}{(}\PY{l+s+s1}{\PYZsq{}}\PY{l+s+s1}{Power of the Test}\PY{l+s+s1}{\PYZsq{}}\PY{p}{)}\PY{p}{;}
\end{Verbatim}
\end{tcolorbox}

    \textbf{(e)} Interpret the graph. What is the test likely to do if the
true value of \(\mu\) is far from 20, and what does the power converge
to (be careful!) when the true value gets close to \(20\)?

    \subsection{Submission Instructions}\label{submission-instructions}

Many assignments throughout the course will have a written portion and a
code portion. Please follow the directions below to properly submit both
portions.

\subsubsection{Written Portion}\label{written-portion}

\begin{itemize}
\tightlist
\item
  Scan all the pages into a PDF. You can use any scanner or a phone
  using applications such as CamScanner. Please \textbf{DO NOT} simply
  take pictures using your phone.
\item
  Please start a new page for each question. If you have already written
  multiple questions on the same page, you can crop the image in
  CamScanner or fold your page over (the old-fashioned way). This helps
  expedite grading.
\item
  It is your responsibility to check that all the work on all the
  scanned pages is legible.
\item
  If you used \(\LaTeX\) to do the written portions, you do not need to
  do any scanning; you can just download the whole notebook as a PDF via
  LaTeX.
\end{itemize}

\subsubsection{Code Portion}\label{code-portion}

\begin{itemize}
\tightlist
\item
  Save your notebook using
  \texttt{File\ \textgreater{}\ Save\ Notebook}.
\item
  Generate a PDF file using
  \texttt{File\ \textgreater{}\ Save\ and\ Export\ Notebook\ As\ \textgreater{}\ PDF}.
  This might take a few seconds and will automatically download a PDF
  version of this notebook.

  \begin{itemize}
  \tightlist
  \item
    If you have issues, please post a follow-up on the general Homework
    7 Ed thread.
  \end{itemize}
\end{itemize}

\subsubsection{Submitting}\label{submitting}

\begin{itemize}
\tightlist
\item
  Combine the PDFs from the written and code portions into one PDF.
  \href{https://smallpdf.com/merge-pdf}{Here} is a useful tool for doing
  so.
\item
  Submit the assignment to Homework 7 on Gradescope.
\item
  \textbf{Make sure to assign each page of your pdf to the correct
  question.}
\item
  \textbf{It is your responsibility to verify that all of your work
  shows up in your final PDF submission.}
\end{itemize}

If you are having difficulties scanning, uploading, or submitting your
work, please read the
\href{https://edstem.org/us/courses/94054/discussion/7533569}{Ed Thread}
on this topic and post a follow-up on the general Homework 7 Ed thread.

    \subsection{\texorpdfstring{\textbf{We will not grade assignments which
do not have pages selected for each
question.}}{We will not grade assignments which do not have pages selected for each question.}}\label{we-will-not-grade-assignments-which-do-not-have-pages-selected-for-each-question.}


    % Add a bibliography block to the postdoc
    
    
   







\begin{align*}
	\Cov(R, B) &= \Cov(\sum_{j=1}^{380}X_{j}, \sum_{k=1}^{380}Y_{k}) \\
	&= \sum_{i=1}^{380}\Cov(X_{i}, Y_i) + 0\\
	&= 380 \cdot \left( E[X_iY_i] - E[X_i]E[Y_i] \right)  \\
	&= 380 \cdot \left( 0 - \frac{18}{38} \cdot \frac{18}{38} \right)  \\
	&= -\frac{1620}{19} \approx -85.26 \\
\end{align*}
\end{document}
